% Proposed additive clarification. Do not replace the pinned manuscript.
% Intended immediately before the two-anchor lemma at Brox lines 2641--2670.
% Covers the inherited precision 16 and 64 tests, and precision 52 in the quartic.

\begin{lemma}[Reciprocal logarithms at every unramified place]
Let $K=\mathbb Q(\sqrt2)$, $D=(3-\sqrt2)^2$, and
$E:y^2=x^3-Dx$. Let $p>76176$ be a rational prime, and let $\mathfrak p$ be any place
of $K$ above $p$. Normalize $v_{\mathfrak p}(p)=1$.
For a nonidentity point $Q$ in the formal group, put
$f=1/x(Q)$ and $z=-x(Q)/y(Q)$. Define
\[
 a_0=1,\qquad a_1=-2/5,\qquad
 a_j=\frac{4j-5}{4j+1}a_{j-1}\quad(j\ge2),\qquad
 L_e(Q)=z\sum_{j=0}^{e/4-1}a_j(Df^2)^j.
\]
For $e\in\{16,52,64\}$ one has
\[
 v_{\mathfrak p}(\log_E Q)=v_{\mathfrak p}(z)=:r\ge1,
 \qquad
 v_{\mathfrak p}(\log_E Q-L_e(Q))\ge(e+1)r,
 \qquad
 \frac{L_e(Q)}{\log_E Q}\in1+p^e\mathcal O_{K_{\mathfrak p}}.
\]
Consequently, for nonidentity formal points $Q,Q'$ and
$e,e'\in\{16,52,64\}$, the truncated ratio has the same valuation
as the full logarithm ratio, and the relative error has valuation
at least $\min(e,e')$. If the full ratio is integral, this is also
an absolute error bound.
\end{lemma}

\begin{proof}
The extension $K_{\mathfrak p}/\mathbb Q_p$ is unramified, whether
$p$ splits or is inert, and $D$ is a unit since $N_{K/\mathbb Q}(D)=49$.
The curve has good reduction. The formal parameter identities are
\[
 f=z^2(1-Df^2),\quad w=Df^2,\quad
 Dz^4=\frac{w}{(1-w)^2},\quad v_{\mathfrak p}(w)=4r.
\]
The recurrence gives
\[
 4w(1-w)G'(w)+(1+w)G(w)=1-w,
 \qquad G(w)=\sum_{j\ge0}a_jw^j.
\]
Differentiating the parameter identity proves that the derivative
of $zG(w)$ is $(1+4Dz^4)^{-1/2}$ with constant term one. Thus
$zG(w)$ is the normalized formal logarithm.

The denominator of $a_j$ divides
$5\prod_{k=2}^j(4k+1)$, hence divides $(4j+1)!$. It follows that
\[
 v_{\mathfrak p}(a_j)\ge-v_p((4j+1)!)
 \ge-\frac{4j+1}{p-1}.
\]
Write $m=e/4\le16$. For $j\ge m$ with $4j+1<p$,
the coefficient is integral and $v(a_jw^j)\ge4mr$.
Otherwise $j\ge19044$, so $3j\ge4m$ and
$(4j+1)/(p-1)\le j$. Therefore
\[
 v(a_jw^j)\ge4jr-j\ge3jr\ge4mr.
\]
The lower bounds tend to infinity. Multiplication by $z$ gives
the asserted error bound. The leading term of the full logarithm
has valuation $r$, and all later terms have larger valuation;
equivalently one may use formal-logarithm isometry on the formal
kernel. Division gives the relative estimate. Dividing the two
relative estimates proves the ratio assertions.
No step uses that the residue field has degree one.
\end{proof}

\begin{corollary}[Two anchors, including inert primes]
Assume the cyclic-group input. Let $H=[h]P$, $U_i=[u_i]P$,
$Q_i=[n_i]P$, $V_0=H-[4]P$, and $V_i=U_i-[4]Q_i$.
Suppose the point and denominator guards of the shared construction
hold and the order-one tests put $H,U_i$ in the formal group at
$\mathfrak p$ and $V_0,V_i$ in the formal group at
$\mathfrak p^\sigma$. Let $k_i,e'\in\{16,52,64\}$ with $e'\ge k_i$,
use numerator precision $k_i$ and denominator precision $e'$, and suppose
$v_{\mathfrak p}(L_i-L'_i)\ge k_i$.
Then $L_i$ is integral and
\[
 L_i\equiv n_i\pmod{p^{k_i}\mathcal O_{K_{\mathfrak p}}}.
\]
\end{corollary}

\begin{proof}
The full logarithm ratios, after conjugating the second one, are
$r_i=u_i/h$ and $r'_i=(u_i-4n_i)/(h-4)$. At least one of $h,h-4$
is a $p$-unit. Its full ratio is integral, so its truncated ratio
is integral by the lemma. The difference test makes the other
truncated ratio integral. Valuation equality then makes both full
ratios integral, and the relative errors are absolute errors.
Hence $L_i\equiv r_i\equiv r'_i\pmod{p^{k_i}}$. Since
\[
 h r_i-(h-4)r'_i=4n_i,
\]
we get $4L_i\equiv4n_i\pmod{p^{k_i}}$ and can divide by four.
If $p$ is inert, conjugation acts within the same unramified
completion; the calculation is unchanged.
\end{proof}

% This clarification adds no variables, guards, exceptional primes, or
% arithmetic hypothesis. It does not certify the missing coefficient DAG
% or its reported degree/height majorants.
